\documentclass{article}%
\usepackage[T1]{fontenc}%
\usepackage[utf8]{inputenc}%
\usepackage{lmodern}%
\usepackage{textcomp}%
\usepackage{lastpage}%
\usepackage{geometry}%
\geometry{margin=1in,headheight=14pt,headsep=25pt}%
\usepackage{hyperref}%
\usepackage{enumitem}%
\usepackage{fancyhdr}%
\usepackage{xcolor}%
\usepackage{url}%
\usepackage{breakurl}%
\usepackage{amsmath}%
\usepackage{amssymb}%
\usepackage{mathtools}%
\usepackage{amsthm}%
\usepackage{thmtools}%
%

            \hypersetup{
                colorlinks=true,
                linkcolor=blue,
                filecolor=magenta,
                urlcolor=blue,
                breaklinks=true,
                bookmarks=true,
                pdfborder={0 0 0}
            }
            \urlstyle{same}
            \Urlmuskip=0mu plus 1mu  % URL breaking in nicer spots
            
            % Math configuration
            \everymath{\displaystyle}
            \setlength{\jot}{10pt}  % Increase spacing between equations
            
            \pagestyle{fancy}
            \fancyhf{}
            \rhead{Generated on \today}
            \lhead{Course Summary}
            \cfoot{\thepage}
            
            % Better Unicode support
            \usepackage[utf8]{inputenc}
            \usepackage{newunicodechar}
            \usepackage[breaklinks=true]{hyperref}
            \setlength{\itemsep}{1em} % Adds spacing between items
            
            % Configure itemize settings globally
            \setlist[itemize]{nosep}
            \setlist[itemize]{leftmargin=*}
        %
\title{Linear Programming Course Summary}%
\author{Generated by Scribe.AI}%
\date{\today}%
%
\begin{document}%
\normalsize%
\maketitle%
\section{Key Terms}%
\label{sec:KeyTerms}%
\subsection{Linear Programming Fundamentals}%
\label{subsec:LinearProgrammingFundamentals}%
\begin{itemize}[leftmargin=*]%
\item \href{https://www.math.purdue.edu/~yipn/421/2024-08-27-ExSimplex.pdf#page=1}{\textbf{Slack Variable}}: A variable added to a less-than-or-equal-to inequality constraint to transform it into an equality constraint.  The slack variable represents the difference between the left-hand side and the right-hand side of the inequality.%
\item \href{https://www.math.purdue.edu/~yipn/421/2024-08-27-ExSimplex.pdf#page=4}{\textbf{Feasible Region}}: The set of all points that satisfy all the constraints of a linear programming problem.%
\item \href{https://www.math.purdue.edu/~yipn/421/2024-08-27-ExSimplex.pdf#page=5}{\textbf{Vertex}}: A point where two or more constraint boundaries intersect in the feasible region of a linear programming problem.  In the Simplex method, optimal solutions are found at vertices.%
\item \href{https://www.math.purdue.edu/~yipn/421/2024-08-27-ExSimplex.pdf#page=6}{\textbf{Basic Variable}}: In the Simplex method, a variable that is expressed in terms of non-basic variables in a dictionary.  Basic variables are set to values that satisfy the constraints.%
\item \href{https://www.math.purdue.edu/~yipn/421/2024-08-27-ExSimplex.pdf#page=6}{\textbf{Non-basic Variable}}: In the Simplex method, a variable that is set to zero in a given iteration of the algorithm.  The values of the basic variables are determined by the values of the non-basic variables.%
\end{itemize}

%
\subsection{Simplex Method Mechanics}%
\label{subsec:SimplexMethodMechanics}%
\begin{itemize}[leftmargin=*]%
\item \href{https://www.math.purdue.edu/~yipn/421/2024-08-27-ExSimplex.pdf#page=9}{\textbf{Pivot Operation}}: The process of exchanging a basic variable with a non-basic variable in the Simplex method, leading to a new dictionary and a movement to a different vertex in the feasible region.%
\item \href{https://www.math.purdue.edu/~yipn/421/2024-08-27-ExSimplex.pdf#page=26}{\textbf{Dictionary}}: A representation of a linear programming problem in a tabular format, where basic variables are expressed as functions of non-basic variables.  It is used in the Simplex method to systematically move from one vertex to another.%
\end{itemize}

%
\subsection{Auxiliary Problem and Initialization}%
\label{subsec:AuxiliaryProblemandInitialization}%
\begin{itemize}[leftmargin=*]%
\item \href{https://www.math.purdue.edu/~yipn/421/2024-08-27-ExSimplex.pdf#page=32}{\textbf{Auxiliary Problem}}: A modified linear programming problem introduced to find an initial feasible solution when the origin is not feasible in the original problem.  It involves adding a new variable to shift the feasible region.%
\end{itemize}

%
\section{Problem Types}%
\label{sec:ProblemTypes}%
\subsection{Linear Programming Fundamentals}%
\label{subsec:LinearProgrammingFundamentals}%
\begin{itemize}[leftmargin=*]%
\item \href{https://www.math.purdue.edu/~yipn/421/2024-08-27-ExSimplex.pdf#page=2}{\textbf{Solving Linear Programs Graphically in Two Dimensions}}: This problem focuses on solving linear programming problems with two variables by visualizing the feasible region and identifying optimal solutions at vertices.%
\end{itemize}

%
\subsection{Simplex Method Techniques}%
\label{subsec:SimplexMethodTechniques}%
\begin{itemize}[leftmargin=*]%
\item \href{https://www.math.purdue.edu/~yipn/421/2024-08-27-ExSimplex.pdf#page=6}{\textbf{Solving Linear Programs Using the Simplex Method}}: This problem involves applying the Simplex algorithm to iteratively find an optimal solution by moving between vertices of the feasible region, updating a dictionary of variables in each iteration.%
\item \href{https://www.math.purdue.edu/~yipn/421/2024-08-27-ExSimplex.pdf#page=26}{\textbf{Understanding the Simplex Dictionary}}: This problem involves understanding the structure and interpretation of the Simplex dictionary, which represents the basic and non-basic variables and their relationships in a linear program.%
\end{itemize}

%
\subsection{Handling Infeasibility and Optimality}%
\label{subsec:HandlingInfeasibilityandOptimality}%
\begin{itemize}[leftmargin=*]%
\item \href{https://www.math.purdue.edu/~yipn/421/2024-08-27-ExSimplex.pdf#page=30}{\textbf{Finding an Initial Feasible Solution}}: This problem addresses situations where the origin is not a feasible solution to the linear programming problem, requiring a method to find a starting feasible point for algorithms like the Simplex method.%
\item \href{https://www.math.purdue.edu/~yipn/421/2024-08-27-ExSimplex.pdf#page=30}{\textbf{Handling Infeasible Linear Programs}}: This problem involves determining whether a given linear program has any feasible solutions at all, and if not, providing a method to determine this.%
\end{itemize}

%
\section{Algorithm Solutions}%
\label{sec:AlgorithmSolutions}%
\subsection{Linear Optimization Algorithms}%
\label{subsec:LinearOptimizationAlgorithms}%
\begin{itemize}[leftmargin=*]%
\item \href{https://www.math.purdue.edu/~yipn/421/2024-08-27-ExSimplex.pdf#page=6}{\textbf{Simplex Method}}: Iteratively move from one vertex of the feasible region to another by setting non-basic variables to zero and choosing a new set of (non)basic variables to improve the objective function value. This corresponds to moving from vertex to vertex in the feasible region.%
\item \href{https://www.math.purdue.edu/~yipn/421/2024-08-27-ExSimplex.pdf#page=32}{\textbf{Auxiliary Problem}}: Introduce a new variable $x_0$ and modify constraints to create a new problem where the origin is feasible. Solve this auxiliary problem to find a feasible solution for the original problem. If $x_0 = 0$ in the optimal solution of the auxiliary problem, then the solution is feasible for the original problem.%
\end{itemize}

%
\end{document}